\section{民事权利}
\begin{enumerate}
  \item 民法是权利法，民法体系的构建以民事权利为逻辑起点和核心。民事权利，是指民事主体为实现某种利益而\zd{依法为或不为某种行为的自由}。
  \item 法定分类。
        \begin{enumerate}
          \item \zd{财产权}。物权、债权。\zd{物债两分，区分原则}。
          \item 综合性权利。知识产权、继承权、股权。
          \item 人身权。人格权、身份权。
        \end{enumerate}
  \item 学理分类。
        \begin{enumerate}
          \item 支配权/管领权。对权利的客体\zd{直接支配}并\zd{排除他人干涉的权利}。
                \begin{enumerate}
                  \item 种类。客体的不同。物权、人身权、知识产权。
                  \item 特征。
                        \begin{enumerate}
                          \item 对世性。权利主体特定而义务主体不特定。
                          \item 特定性。\zd{客体是特定化的}财产、人身利益和无形财产。
                          \item 直接性。权利的行使和实现无需义务人的积极作为，义务人只负有消极不作为义务。
                          \item 排他性。\zd{不容在同一标的上设定不相容的第二个权利}。
                        \end{enumerate}
                \end{enumerate}
          \item \zd{形成权}。依权利人单方（一方）意思表示而使民事法律关系产生、变更或消灭（即“私法效果”）的权利。
                \begin{enumerate}
                  \item 产生。当事人约定 \hz 法律规定。
                  \item 行使。
                        \begin{enumerate}
                          \item 明示和默示。
                          \item 行使是否需要特定形式。简单形成权/单纯形成权 \yjt 单方意思表示即可发生效力；形成诉权。
                        \end{enumerate}
                  \item 9 类典型\blh。撤销权、（法定）抵销权、追认权、解除权、否认权、选择权、继承人沉默接受继承、受遗赠人沉默放弃受遗赠。
                  \item 期间限制。除斥期间而非诉讼时效的限制。
                \end{enumerate}
          \item 请求权。特定的权利人\zd{请求特定的义务人为或不为特定行为的权利}。进攻性，矛。
                \begin{enumerate}
                  \item 特征。
                        \begin{enumerate}
                          \item 相对性（权利人与义务人均特定，\zd{效力仅及于当事人之间}），权利人实现其权利必须具体特定人协助。
                          \item 客体为特定行为，包括作为和不作为。
                          \item 非公示性。
                        \end{enumerate}
                  \item 发生。既可以基于基础权利\zd{受侵害}而发生，亦可以基于基础权利\zd{未受侵害}而发生。
                \end{enumerate}
          \item 抗辩权。一时性或永久对抗对方请求权的权利。防御性，盾。
                \begin{enumerate}
                  \item 广义上的抗辩权是妨碍他人行使其权利、提出免除或减轻其民事责任主张的对抗权。
                  \item 狭义的抗辩权则是指\zd{专门对抗请求权}的权利。
                  \item 一时性/延期性/延缓性\blh。可以暂时地拒绝相对人的请求权效力。
                        \begin{enumerate}
                          \item 双务合同。同时履行抗辩权、先履行抗辩权、不安抗辩权。
                          \item 一般保证人的先诉抗辩权。
                        \end{enumerate}
                  \item 永久性/消灭时效\blh。
                  \item 特征。
                        \begin{enumerate}
                          \item 以请求权的行使且承认对方有请求权为前提。
                          \item 有期间限制。
                          \item \zd{不主张视为放弃，法院一般不得依其职权主动审查，也不得主动释明}。
                        \end{enumerate}
                \end{enumerate}
        \end{enumerate}
  \item 民事权利的行使和限制。以鼓励和维护为原则，以限制或否定为例外。
\end{enumerate}
